% GENERATED FILE. Edit canonical lessons, then run npm run generate:books.
% canonical-source-hash: fnv1a64:78036a8e65efc962
% canonical-lessons: FA-C532-ekran, FA-C532-kontorol, FA-C532-safhe-ye-namayesh, FA-C532-padkast, FA-C532-hedfon

\chapter{Watching and Listening}
\label{ch:fa-watching-and-listening}
\hlchaptermodality{532}

\begin{chapteropening}
\textbf{By the end of this chapter:} \emph{I can talk about screens and listening.}

Complete the last lesson of chapter 532: I can talk about screens and listening.
\end{chapteropening}

\section[\texorpdfstring{\fa{اکران} (ekrân)}{اکران (ekrân)}]{\fa{اکران} (ekrân) — a screening, a (film) release}

\label{lesson:FA-C532-ekran}



\begin{quote}
Before the new one: say the Persian for dubbed, then the Persian for a TV channel.
\end{quote}

\subsection*{You'll want to know: \fa{اکران}}
\textbf{\fa{اکران}} (\emph{ekrân}) — \textquotedblleft{}a screening, a (film) release\textquotedblright{}.

\begin{quote}
\fa{فیلم} \fa{تازه} \fa{او} \fa{از} \fa{هفته} \fa{بعد} \fa{اکران} \fa{میشود}.
\end{quote}

\begin{quote}
\emph{film-e tâze-ye u az hafte-ye ba'd ekrân mishavad.}
\end{quote}

\begin{quote}
His new film opens next week.
\end{quote}

\emph{ekrân shodan} — to be screened.

\begin{cousinweb}
\emph{ekrân} is a loan, usually traced to French \emph{écran}, 'screen'; Persian took the screen and made it the screening. EtymDB has no entry for it.
\end{cousinweb}

\begin{grammarlens}[title={What you've built}]
One more word for telling, explaining and describing at B1.
\end{grammarlens}

\subsection*{Your turn}
\begin{itemize}
\raggedright
  \item \emph{Say it:} \emph{ekrân}
  \item \emph{Say it:} \emph{ekrân}, once more
  \item \emph{Say it:} a sentence of your own with \emph{ekrân}
  \item \emph{Recall:} say the Persian for dubbed, then the Persian for a TV channel, then say \emph{ekrân} again
\end{itemize}

\begin{tcolorbox}[breakable,colback=teal!4,colframe=teal!35!black,title={Before you move on}]
What does \fa{اکران} mean? (\textquotedblleft{}A screening, a (film) release\textquotedblright{}.) Say it once more.
\end{tcolorbox}
\section[\texorpdfstring{\fa{کنترل} (kontorol)}{کنترل (kontorol)}]{\fa{کنترل} (kontorol) — a remote control; control}

\label{lesson:FA-C532-kontorol}



\begin{quote}
Before the new one: say the Persian for a TV channel, then the Persian for a screening.
\end{quote}

\subsection*{You'll want to know: \fa{کنترل}}
\textbf{\fa{کنترل}} (\emph{kontorol}) — \textquotedblleft{}a remote control; control\textquotedblright{}.

\begin{quote}
\fa{کنترل} \fa{تلویزیون} \fa{زیر} \fa{مبل} \fa{بود}.
\end{quote}

\begin{quote}
\emph{kontorol-e televiziyon zir-e mobl bud.}
\end{quote}

\begin{quote}
The TV remote was under the sofa.
\end{quote}

\emph{kontorol kardan} — to check, to control.

\begin{cousinweb}
\emph{kontorol} is a modern loan, usually traced to French \emph{contrôle}. EtymDB has no entry for the Persian word.
\end{cousinweb}

\begin{grammarlens}[title={What you've built}]
One more word for telling, explaining and describing at B1.
\end{grammarlens}

\subsection*{Your turn}
\begin{itemize}
\raggedright
  \item \emph{Say it:} \emph{kontorol}
  \item \emph{Say it:} \emph{kontorol}, once more
  \item \emph{Say it:} a sentence of your own with \emph{kontorol}
  \item \emph{Recall:} say the Persian for a TV channel, then the Persian for a screening, then say \emph{kontorol} again
\end{itemize}

\begin{tcolorbox}[breakable,colback=teal!4,colframe=teal!35!black,title={Before you move on}]
What does \fa{کنترل} mean? (\textquotedblleft{}A remote control; control\textquotedblright{}.) Say it once more.
\end{tcolorbox}
\section[\texorpdfstring{\fa{صفحه} \fa{نمایش} (safhe-ye namâyesh)}{صفحه نمایش (safhe-ye namâyesh)}]{\fa{صفحه} \fa{نمایش} (safhe-ye namâyesh) — a screen (display)}

\label{lesson:FA-C532-safhe-ye-namayesh}



\begin{quote}
Before the new one: say the Persian for a screening, then the Persian for a remote control.
\end{quote}

\subsection*{You'll want to know: \fa{صفحه} \fa{نمایش}}
\textbf{\fa{صفحه} \fa{نمایش}} (\emph{safhe-ye namâyesh}) — \textquotedblleft{}a screen (display)\textquotedblright{}.

\begin{quote}
\fa{صفحه} \fa{نمایش} \fa{موبایلم} \fa{شکست}.
\end{quote}

\begin{quote}
\emph{safhe-ye namâyesh-e mobâylam shekast.}
\end{quote}

\begin{quote}
My phone screen broke.
\end{quote}

\emph{safhe} 'page, surface' + \emph{namâyesh} 'showing'.

\begin{cousinweb}
\emph{safhe-ye namâyesh} is built in Persian on \emph{safhe}, 'page, surface', from Arabic \emph{ṣafḥa} (EtymDB), and native \emph{namâyesh}, 'showing', from \emph{namudan}.
\end{cousinweb}

\begin{grammarlens}[title={What you've built}]
One more word for telling, explaining and describing at B1.
\end{grammarlens}

\subsection*{Your turn}
\begin{itemize}
\raggedright
  \item \emph{Say it:} \emph{safhe-ye namâyesh}
  \item \emph{Say it:} \emph{safhe-ye namâyesh}, once more
  \item \emph{Say it:} a sentence of your own with \emph{safhe-ye namâyesh}
  \item \emph{Recall:} say the Persian for a screening, then the Persian for a remote control, then say \emph{safhe-ye namâyesh} again
\end{itemize}

\begin{tcolorbox}[breakable,colback=teal!4,colframe=teal!35!black,title={Before you move on}]
What does \fa{صفحه} \fa{نمایش} mean? (\textquotedblleft{}A screen (display)\textquotedblright{}.) Say it once more.
\end{tcolorbox}
\section[\texorpdfstring{\fa{پادکست} (pâdkast)}{پادکست (pâdkast)}]{\fa{پادکست} (pâdkast) — a podcast}

\label{lesson:FA-C532-padkast}



\begin{quote}
Before the new one: say the Persian for a remote control, then the Persian for a screen.
\end{quote}

\subsection*{You'll want to know: \fa{پادکست}}
\textbf{\fa{پادکست}} (\emph{pâdkast}) — \textquotedblleft{}a podcast\textquotedblright{}.

\begin{quote}
\fa{او} \fa{در} \fa{راه} \fa{کار} \fa{پادکست} \fa{گوش} \fa{میدهد}.
\end{quote}

\begin{quote}
\emph{u dar râh-e kâr pâdkast gush midahad.}
\end{quote}

\begin{quote}
He listens to podcasts on the way to work.
\end{quote}

Plural \emph{pâdkasthâ}.

\begin{cousinweb}
\emph{pâdkast} was borrowed from English \emph{podcast} (EtymDB).
\end{cousinweb}

\begin{grammarlens}[title={What you've built}]
One more word for telling, explaining and describing at B1.
\end{grammarlens}

\subsection*{Your turn}
\begin{itemize}
\raggedright
  \item \emph{Say it:} \emph{pâdkast}
  \item \emph{Say it:} \emph{pâdkast}, once more
  \item \emph{Say it:} a sentence of your own with \emph{pâdkast}
  \item \emph{Recall:} say the Persian for a remote control, then the Persian for a screen, then say \emph{pâdkast} again
\end{itemize}

\begin{tcolorbox}[breakable,colback=teal!4,colframe=teal!35!black,title={Before you move on}]
What does \fa{پادکست} mean? (\textquotedblleft{}A podcast\textquotedblright{}.) Say it once more.
\end{tcolorbox}
\section[\texorpdfstring{\fa{هدفون} (hedfon)}{هدفون (hedfon)}]{\fa{هدفون} (hedfon) — headphones}

\label{lesson:FA-C532-hedfon}



\begin{quote}
Before the new one: say the Persian for a screen, then the Persian for a podcast.
\end{quote}

\subsection*{You'll want to know: \fa{هدفون}}
\textbf{\fa{هدفون}} (\emph{hedfon}) — \textquotedblleft{}headphones\textquotedblright{}.

\begin{quote}
\fa{او} \fa{با} \fa{هدفون} \fa{موسیقی} \fa{گوش} \fa{میدهد}.
\end{quote}

\begin{quote}
\emph{u bâ hedfon musiqi gush midahad.}
\end{quote}

\begin{quote}
She listens to music on headphones.
\end{quote}

One word for the pair.

\begin{cousinweb}
\emph{hedfon} is a modern loan, usually traced to English \emph{headphone}. EtymDB has no entry for it.
\end{cousinweb}

\begin{grammarlens}[title={What you've built}]
One more word for telling, explaining and describing at B1.
\end{grammarlens}

\subsection*{Your turn}
\begin{itemize}
\raggedright
  \item \emph{Say it:} \emph{hedfon}
  \item \emph{Say it:} \emph{hedfon}, once more
  \item \emph{Say it:} a sentence of your own with \emph{hedfon}
  \item \emph{Recall:} say the Persian for a screen, then the Persian for a podcast, then say \emph{hedfon} again
\end{itemize}

\begin{tcolorbox}[breakable,colback=teal!4,colframe=teal!35!black,title={Before you move on}]
What does \fa{هدفون} mean? (\textquotedblleft{}Headphones\textquotedblright{}.) Say it once more.
\end{tcolorbox}
